% 第 6 题三视图：按固定源 PDF 的投影轮廓与尺寸线位置重绘（尺寸数字只保留原图印出的 1、1、1、1）。
\begin{tikzpicture}[every node/.style={font=\small},scale=0.9,line cap=round,line join=round]
% 正（主）视图：尖顶在左上，底边在右下，无尺寸标注
\draw (0,1.0)--(0.98,0)--(1.98,0)--cycle;
\node[below] at (0.99,-0.28) {正（主）视图};
% 侧（左）视图：直角在左下；左侧竖直尺寸线 1，下方水平尺寸线 1
\draw (2.9,0)--(2.9,1.0)--(3.88,0)--cycle;
\draw[<->] (2.72,0)--(2.72,1.0); \node[left] at (2.72,0.5) {1};
\draw[<->] (2.9,-0.20)--(3.88,-0.20); \node[below] at (3.39,-0.20) {1};
\node[below] at (3.39,-0.62) {侧（左）视图};
% 俯视图：直角在右上；上边长 2（两段各 1），上边中点到右下顶点为虚线
\draw (0,-1.98)--(1.98,-1.98)--(1.98,-2.96)--cycle;
\draw[dashed] (0.99,-1.98)--(1.98,-2.96);
\draw[<->] (0,-1.76)--(0.99,-1.76); \node[above] at (0.495,-1.76) {1};
\draw[<->] (0.99,-1.76)--(1.98,-1.76); \node[above] at (1.485,-1.76) {1};
\node[below] at (0.99,-3.30) {俯视图};
\end{tikzpicture}
